
\chapter{Results}
%%%%%%%%%%%%%%%%%%%%%%%%%%%%%%%%%%%%%%%%%%%%%%%%%%%%%%%%%%%%%%%%%%%%%%%%%%%%%%%%%%%
\section{Expression and localization of eGFP-uTP constructs in \textit{Angomonas deanei}}
%---------------------------------------------------------------------------------------------------
\subsection{Generation of strains expressing eGFP-uTP constructs}
For the expression of the nitroplast targeting signals in \glsxtrshort{Adea}, they needed to be integrated into the genome. The strain Adea126 ($\Delta$\textit{\gls{g-ama}\textsuperscript{hygR, msc-etp1}}; see Table~\ref{tab:Adea_list}) was chosen to be transfected with these \gls{egfp}-\gls{utp} constructs for possible co-localization. Adea126 contained \glsxtrlong{mSc}-\gls{etp}1 in the \gls{g-ama} locus; therefore, three plasmids had to be cloned that contained the transfection cassette for the \gls{d-ama} locus with \gls{neoR} complementing \gls{hygR} of the strain, an \gls{egfp}-tag for verification, and each one of the three \gls{utp}s. The three plasmids were successfully cloned using Gibson assembly and were named pAdea481, pAdea482, and pAdea483, which had \gls{utp_B}, \gls{utp_E} and \gls{utp_C}, respectively (see Table~\ref{tab:plasmid_list}).  

\par The strain Adea126, hereafter also referred to as the background strain, was transfected with the generated cassettes. To verify the genomic integration of the transfection cassette with the \gls{egfp}-tagged \gls{utp}s, a \gls{TDPCR} was performed. The amplification of the \gls{d-ama} locus and its surrounding genomic region using primers 49 and 545 yielded fragments with the sizes of 3.3 kbp for \gls{WT} and 5.2 kbp for the transfectants that contained \gls{egfp}-\gls{utp}\textsubscript{X}. Two clones (c1 and c2) of each construct, which were heterozygous for the \gls{d-ama} locus containing one allele with the transfection cassette and another allele for \gls{WT}, were maintained as continuous cultures (see Figure~\ref{fig:TDPCR_Ango}). 

\begin{figure}[H]
        \centering
        \includegraphics[width=15cm]{"4.1.1_TDPCR_Ango_PfkB_PyrE"}
\caption{
\textbf{Verification of the genomic integration in transfected Adea126 clones}
    \gls{TDPCR} products from extracted \gls{gDNA} were loaded onto a 1~\% agarose gel. Primers 49 and 545 were used; GeneRuler 1 kbp (L); Wildtype (WT); c1--c6 being clones number 1 to 6. pAdea483 containing the \gls{utp_B}. pAdea482 containing the \gls{utp_E}.} 
		\label{fig:TDPCR_Ango}
\end{figure}

\par The generated strains will be referred to as Adea538 ($\Delta$\textit{\gls{g-ama}\textsuperscript{hygR, msc-etp1}}, $\Delta$\textit{\gls{d-ama}\textsuperscript{neoR, egfp-utp\textsubscript{pfkB}}}) and Adea539 ($\Delta$\textit{\gls{g-ama}\textsuperscript{hygR, msc-etp1}}, $\Delta$\textit{\gls{d-ama}\textsuperscript{neoR, egfp-utp\textsubscript{pyrE}}}). The DNA sequences that were specifically transferred are listed in Table~\ref{tab:Adea_list}. After several attempts, it was not possible to transfect Adea126 with the transfection cassette containing the sequence for \gls{utp_C} (pAdea481). Thus, the work regarding heterologous expression in \glsxtrshort{Adea} focuses only on uTP\textsubscript{PfkB} and uTP\textsubscript{PyrE}.

%---------------------------------------------------------------------------------------------------
%---------------------------------------------------------------------------------------------------
\subsection{Heterologous expression and fluorescence assay of eGFP-tagged uTPs}
After the successful integration of the nitroplast transit peptides tagged with \gls{egfp}, the expression was then analyzed via Western blotting and in-gel fluorescence assay to verify expression and fluorescence of the constructs. 

\par First, a Western blot analysis was done to verify the expression of \gls{egfp}-tagged uTPs in Adea538 and Adea539. Bands were detected in the sample of Adea304 ($\Delta$\textit{\gls{g-ama}\textsuperscript{neoR, egfp}}; see Table~\ref{tab:Adea_list}) at around 30~kDa and in the samples of the two clones of Adea538 and Adea539 at around 55~kDa (see Figure~\ref{fig:Western_Ango}B). The uTP fusion proteins ran consistently higher (around 55 kDa) than calculated by ProtParam (46~kDa for \gls{egfp}-\gls{utp_B} and 48~kDa for \gls{egfp}-\gls{utp_E}). Moreover, a comparison of the fainter bands from the Adea538 and Adea539 samples to the bright band from the Adea304 sample indicates a difference in expression levels, which is not due to loading of different amounts of protein sample (see Figure~\ref{fig:Western_Ango}A). However, the \gls{egfp}-\gls{utp} constructs had similar expression levels to each other. \glsxtrshort{Adea} can generally produce the \gls{egfp}-\gls{utp} constructs. Though the weak expression might cause problems for later analysis  

\begin{figure}[H]
        \centering
        \includegraphics[width=15cm]{"4.1.2_TCE_And_Western_Ango"}
\caption{
\textbf{\gls{page} and Western blot of the expression of \gls{egfp}-\gls{utp}s from the \gls{d-ama} locus}. 
    Samples from \glsxtrshort{Adea} expression cultures were run on an \gls{page}. Additionally, a Western blot analysis with \textgreek{α}-GFP was done. Prestained PageRuler (L); Adea304 expressed only \gls{egfp}; Adea126 expressed \gls{mSc}-ETP1; Adea539 expressed mSc-ETP1/eGFP-uTP\textsubscript{PfkB}; Adea538 expressed mSc-ETP1/eGFP-uTP\textsubscript{PyrE}.
    (A) Stain-free image of the \gls{page}.
    (B) Corresponding Western blot.}
		\label{fig:Western_Ango}
\end{figure}

\begin{figure}[H]
        \centering
        \includegraphics[width=10cm]{"4.1.2_Ingel_Ango_msc_gfp"}
\caption{
\textbf{In-gel fluoresence assay of semi-denatured \gls{egfp}-\gls{utp} constructs}. 
     Samples from \glsxtrshort{Adea} expression cultures were run on \gls{page}. Prestained PageRuler (L); Wildtype (WT); Adea126 expressed \gls{mSc}-\gls{etp}1; Adea539 expressed \gls{mSc}-\gls{etp}1/\gls{egfp}-\gls{utp_B}; Adea538 expressed \gls{mSc}-\gls{etp}1/\gls{egfp}-\gls{utp_E}; Adea304 expressed only \gls{egfp} and had been loaded onto the \gls{page} with half of the calculated protein amount.
    (A) Coomassie merged with Cy3 channel for \glsxtrlong{mSc} detection.
    (B) Coomassie merged with Cy2 channel for \gls{egfp} detection. For the corresponding \gls{TCE} image merged with the coomassie channel, see Supplementary Figure \ref{sup:ingel_tce}}
       \label{fig:ingel_Ango}
\end{figure}

\par Next, the fluorescence of \gls{egfp} fused to the nitroplast transit peptides was examined by doing an in-gel fluorescence assay (see Figure~\ref{fig:ingel_Ango}). The samples from the expression cultures were semi-denatured and ran on an \gls{page}. The sample for Adea304 was adjusted to be loaded onto the \gls{page} at half the protein amount of the other samples. This was done to minimize the possibility that strong signals from the Adea304 sample prevented the detection of the signals from the \gls{egfp}-tagged \gls{utp}s. There were distinct signals in the \gls{page} for \glsxtrlong{mSc} at around 77~kDa (see Figure~\ref{fig:ingel_Ango}A), which was the \gls{ES} marker (\glsxtrlong{mSc}-\gls{etp}1). Furthermore, signals for \gls{egfp} at around 35~kDa were also detected (Figure \ref{fig:ingel_Ango}B). The positive control and nearly all clones of both strains expressing \gls{egfp}-\gls{utp}\textsubscript{X}, except for Adea539 c1, displayed fluorescence. The brightness of the signals for the \gls{egfp}-tagged \gls{utp}s seemed similar or less bright than the sample of Adea304, which contained half the protein amount. The \gls{utp} fusion proteins displayed fluorescence, indicating that \gls{egfp} adopts native folding even when the \gls{utp}s were fused to its C-terminus and that likely no interaction of the \gls{utp} with \gls{egfp} occurred.  

%The signals for \gls{egfp} in the in-gel fluorescence assay were separately verified by a Western blot analysis (data not shown).



%---------------------------------------------------------------------------------------------------
%---------------------------------------------------------------------------------------------------
\subsection{Localization in \textit{Angomonas deanei} via fluorescence microscopy}
After confirming the expression and fluorescence of \gls{egfp}-\gls{utp} constructs, the \glsxtrshort{Adea} strains expressing these fusion proteins were compared with the background strain and the positive control using epifluorescence (see Figure \ref{fig:epiflu_Ango}) and confocal fluorescence microscopy (see Figure \ref{fig:confo_Ango}) to localize the expressed recombinant nitroplast transit peptides. 

\par The strains were first examined using epifluorescence microscopy. For the overview images of the single cells shown in Figure~\ref{fig:epiflu_Ango}, see Supplementary Figure~\ref{sup:epiflu_Ango}. Adea126 showed fluorescence for \glsxtrlong{mSc} to which \gls{etp}1 is fused; the whole construct localizes around the \gls{ES} as expected \parencite{morales_development_2016}. Adea304 also displayed its characteristic fluorescence. The \gls{egfp} signal was diffused in the cell, indicating cytosolic localization. Moreover, the strains expressing the \gls{utp} fusion proteins showed no abnormal phenotypes, suggesting no disruption of major cellular processes. While the background strain and positive control exhibited typical fluorescence, the \glsxtrshort{Adea} strains expressing the nitroplast transit peptides with \gls{egfp}-tag (Adea538 and Adea539) lacked the signals for \gls{egfp} that exceeded the autofluorescence but showed fluorescence for \glsxtrlong{mSc} and Hoechst (see Figure~\ref{fig:epiflu_Ango}). Furthermore, cells expressing the \gls{utp} construct exhibited no abnormal phenotype in the appearance of either \glsxtrshort{Adea} or its \gls{ES}.

\par To further validate the epifluorescent findings, the same strains but different exponential cultures were analyzed with the Leica TCS SP8 STED 3X. Here, time-gating was utilized to avoid the autofluorescence, and the channels for \gls{dic} and Hoechst were not recorded. Again, Adea126 displayed the typical fluorescence pattern, with \gls{mSc}-\gls{etp}1 localizing around the \gls{ES} (see Figures~\labelcref{fig:epiflu_Ango,fig:confo_Ango}). This was also true for Adea304. The \gls{egfp} signals were distributed throughout the cells, appearing as dots and most likely localizing in the cytosol (see Figures~\labelcref{fig:confo_Ango,fig:epiflu_Ango}). All clones of both strains expressing the eGFP-uTP constructs show fluorescence for \glsxtrlong{mSc} but not for \gls{egfp}, similar to what was observed in the epifluorescence microscopy (see Figures~\labelcref{fig:confo_Ango,fig:epiflu_Ango}). In some images, like in Adea538 c2, there were weak signals for \gls{egfp} that could be attributed to the artifacts of image processing (see Figure~\ref{fig:epiflu_Ango}A).

\par Moreover, proteins precipitated from the supernatant of the harvested expression cultures were also examined by Western blot analysis to see if export occurs. The only band present was for Adea304 at around 35~kDa while the other strains lacked any signals (see Supplementary Figure~\ref{sup:sn_west}). The band for Adea304 might be most likely artifacts of cell remnants inside the sample, though a putative protein export of \gls{egfp} cannot be excluded. 

\par The \gls{egfp}-tagged nitroplast targeting signals should be localizing inside the \glsxtrshort{Adea} cells, as they were not detected outside the cells. Further, \gls{egfp} should be detectable as seen by the in-gel fluorescence assay. Possibly low expression, as seen in the Western blot analysis, leading to low concentration of the fusion proteins, might hinder localization. 

\begin{figure}[H]
        \centering
        \includegraphics[width=16cm]{"4.1.3_All_Closeups_epi"}
\caption{
\textbf{Epifluorescence microscopy images of fixed single \glsxtrshort{Adea} cells}. 
    \glsxtrshort{Adea} strains were grown to the exponential phase when the cells were fixed and stained. The channels are \glsxtrfull{dic}, eGFP, \glsxtrlong{mSc}, and Hoechst.
    Adea126 expressed \gls{mSc}-\gls{etp}1; Adea304 expressed \gls{egfp}; Adea538 expressed \gls{mSc}-\gls{etp}1/\gls{egfp}-\gls{utp_E}; Adea538 expressed \gls{mSc}-\gls{etp}1/\gls{egfp}-\gls{utp_B}.}
       \label{fig:epiflu_Ango}
\end{figure}


\begin{figure}[H]
        \centering
        \includegraphics[width=16cm]{"4.1.3_All_all_confo"}
\caption{
\textbf{Confocal fluorescent microscopy images of fixed \glsxtrshort{Adea} cells}. 
    \glsxtrshort{Adea} strains were grown to the exponential phase when the cells were fixed. The channels were eGFP and \glsxtrlong{mSc}. Adea126 expressed \gls{mSc}-\gls{etp}1; Adea304 expressed \gls{egfp}; Adea538 expressed \gls{mSc}-\gls{etp}1/\gls{egfp}-\gls{utp_E}; Adea538 expressed \gls{mSc}-\gls{etp}1/\gls{egfp}-\gls{utp_B}. (A) Single cells (zoom-in of the overview in percentage) and (B) overview of every strain.}
       \label{fig:confo_Ango}
\end{figure}

%---------------------------------------------------------------------------------------------------
%%%%%%%%%%%%%%%%%%%%%%%%%%%%%%%%%%%%%%%%%%%%%%%%%%%%%%%%%%%%%%%%%%%%%%%%%%%%%%%%%%%
%---------------------------------------------------------------------------------------------------
\section{Heterologous expression of His-tagged uTPs in \textit{Escherichia coli} LOBSTR}
For the functional and structural characterization of the nitroplast targeting signals, it was essential to obtain purified \gls{utp}s in high amounts. This was achieved by overexpressing the affinity-tagged \gls{utp}s  in \glsxtrshort{Eco} and their subsequent purification using \gls{IMAC}. The plasmids for expressing \gls{6xhis}-\gls{sumo}-TEV-3xGS-\gls{utp}\textsubscript{X} (shortened to \gls{6xhis}-\gls{utp}\textsubscript{X} for simplicity) were generated by Victoria Klimenko (see Table \ref{tab:plasmid_list}), who also transformed \glsxtrshort{Eco} LOBSTR with these plasmids (see Table \ref{tab:Eco_list}). An expression and a solubility test were performed to analyze how well EC173 and EC174 expressed His-tagged \gls{utp_C} and \gls{utp_E}, respectively. Furthermore, a purification test was done to determine whether or not the \gls{6xhis}-uTPs can be purified. Due to time constraints, the \gls{utp_B} was not examined here. For each test, an \gls{page} and a corresponding Western blot had been done. 
%---------------------------------------------------------------------------------------------------
\subsection{Expression test: His-tagged uTPs were expressed}
The strains were grown in different media (\gls{lb} and \gls{2yt}) and at different temperatures (37~°C, 28~°C, and 16~°C). Expression of the \gls{6xhis}-\gls{utp}\textsubscript{X} was induced with 0.5 mM \gls{IPTG}.

\begin{figure}[H]
        \centering
        \includegraphics[width=10cm]{"4.2.1_TestSDS_ThrC"}
\caption{
\textbf{\gls{page} of the test expression of the strain EC173 (\gls{utp_C})}. 
     Main cultures were induced with 0.5 mM \gls{IPTG}, and grown in different conditions. Samples were taken accordingly: \glsxtrfull{ni}, 1/2/3/4~hours after induction and \glsxtrfull{on} stain-free imaging of the \gls{page} and prestained PageRuler (L). Cultures were grown in (A) \gls{lb} and (B) \gls{2yt} medium.}
       \label{fig:Test_SDS_C}
\end{figure}

\par For the test expression of the strain EC173, which should produced \gls{6xhis}-\gls{utp_C}, distinctive bands were visible at approximately 40~kDa (see Figure \ref{fig:Test_SDS_C}); running as expected (37~kDa as calculated by Protparam). The expression conditions that yielded the most recombinant protein were: 3~hours after induction at 37~°C and 4~hours after induction at 28°C in both media. Moreover, incubating the induced cultures overnight (see the expression conditions 28~°C 4h and overnight, as well as 16°C overnight in both media; Figure~\ref{fig:Test_SDS_C}A, B) seemed to lead to a decreased amount of the fusion proteins. The Western blot analysis verified that the bands in the \gls{page} were indeed the \gls{6xhis}-\gls{utp}s (see Supplementary Figure~\ref{sup:Test_WB_C}), though additional smaller protein bands (at around 30~kDa) were also detected by the specific antibody in multiple expression conditions. Furthermore, no protein with the \gls{6xhis}-tag was detected in the 16~°C expression conditions in both media, in line with the observations in the \gls{page}s. For further analysis, four conditions were chosen based on the brightness of the bands and that differing growth temperatures: \gls{lb} medium at 37~°C for 3~hours and 28~°C for 4~hours, as well as \gls{2yt} medium at 37~°C for 3~hours and 16~°C overnight.


\begin{figure}[H]
        \centering
        \includegraphics[width=9.5cm]{"4.2.1_TestSDS_PyrE"}
\caption{
\textbf{\gls{page} of the test expression of the strain EC174 (\gls{utp_E})}. 
    Main cultures were induced with 0.5 mM \gls{IPTG} and grown in different conditions. Samples were taken accordingly: \glsxtrfull{ni}, 1/2/3/4~hours after induction and \glsxtrfull{on}; stain-free imaging of the \gls{page} and prestained PageRuler (L). Cultures were grown in (A) \gls{lb} and (B) \gls{2yt} medium.}
       \label{fig:Test_SDS_E}
\end{figure}


\par Similar results for the test expression for the strain EC174, which produced \gls{6xhis}-\gls{utp_E}, were obtained. Bands occurred at around 40~kDa (see Figure \ref{fig:Test_SDS_C}), close to the expected size (36~kDa as calculated by Protparam). The expression conditions that yielded the most recombinant protein were: 3~hours after induction at 37~°C for \gls{lb} media (see Figure~\ref{fig:Test_SDS_E}A), as well as 3~hours after induction at 37~°C and 4~hours after induction at 28°C for \gls{2yt} media (see Figure~\ref{fig:Test_SDS_E}B). Like in the test expression of the strain EC173, incubating the induced cultures overnight does not lead to an increased amount of recombinant protein. Differing to the strain EC173, the expression condition \gls{2yt} 16~°C overnight showed similar amounts of protein compared to the other conditions (see Figure~\ref{fig:Test_SDS_E}). Moreover, the Western blot analysis verified that the bands in the \gls{page} were indeed the \gls{6xhis}-\gls{utp}s (see Supplementary Figure~\ref{sup:Test_WB_E}). For further analysis, the following four conditions were chosen: 37~°C for 3~hours, 28~°C for 4~hours as well as overnight, and 16~°C overnight in \gls{2yt} medium.

%---------------------------------------------------------------------------------------------------
%---------------------------------------------------------------------------------------------------
\subsection{Solubility test: cytosolic fraction contained soluble uTPs}
A solubility test was conducted to obtain an expression condition yielding the most protein in the cytosolic fraction for easier purification. Larger cultures for the four conditions were set up and harvested. Cell lysate samples as well as the supernatant of centrifuged cell lysates were prepared and loaded onto an \gls{page}s for comparison. 

\begin{figure}[H]
        \centering
        \includegraphics[width=10cm]{"4.2.2_SolTest_SDSonly"}
\caption{
\textbf{\gls{page} of the solubility test for the strains EC173 and EC174}. 
    Four conditions were tested. Cell lysate and supernatant from the centrifuged cell lysate were loaded onto the \gls{page}. Stain-free imaging of the \gls{page} with prestained PageRuler (L). 
    (A) EC173 expressing \gls{6xhis}-\gls{utp_C}. 
    (B) EC174 expressing \gls{6xhis}-\gls{utp_E} and grown in \gls{2yt} medium.}
       \label{fig:Soltest_SDSonly}
\end{figure}


\par Both \gls{utp}s displayed prominent bands at around 40~kDa in the cell lysate samples (see Figure~\ref{fig:Soltest_SDSonly}). The 16~°C temperature condition yielded the least amount of recombinant protein in the cell lysate fraction for both \gls{6xhis}-tagged \gls{utp}s, in line with the expression tests. Moreover, a comparison of the cell lysate with the supernatant fraction showed that most of the \gls{6xhis}-\gls{utp} constructs remained in the pellet instead of the cytosolic fraction. For \gls{utp_C}, the expression condition that yielded the most soluble protein in the supernatant compared to its cell lysate was \gls{lb} medium 4~hours after induction at 28~°C (see Figure~\ref{fig:Soltest_SDSonly}A). Similarly, the expression condition \gls{2yt} medium after 4~hours of induction at 28~°C contained the most soluble protein in the supernatant compared to its cell lysate for \gls{utp_E} (see Figure~\ref{fig:Soltest_SDSonly}B). Both \gls{page}s were subjected to Western blot analysis to verify that prominent bands are the \gls{6xhis}-\gls{utp}s, and smaller protein bands at around 30~kDa were again detected  (See Supplementary Figure~\ref{sup:SolTest_Western}).

%---------------------------------------------------------------------------------------------------
%---------------------------------------------------------------------------------------------------
\subsection{Purification test: elutions contained His-tagged uTPs}
For the identification of the optimal procedure for purifying the nitroplast transit peptides with the \gls{6xhis}-tag, a test purificiation was conducted.

\begin{figure}[H]
        \centering
        \includegraphics[width=11cm]{"4.2.3_TestPuri_WesternOnly"}
\caption{
\textbf{Western blot of the test purification for \gls{6xhis}-\gls{utp_C} and \gls{6xhis}-\gls{utp_E}}. 
    Elution was done in steps with increasing concentration of imidazole (50~--~500~mM); prestained PageRuler (L); \glsxtrfull{ft}; first and third centrifugation step (I, III). For corresponding \gls{page} see Supplementary Figure~\ref{sup:TestPuri_SDS}.
    (A) \gls{6xhis}-\gls{utp_C} expressed by EC173, which was grown in \gls{lb} at 28°C and harvested 4 hours after induction as its expression condition. 
    (B) \gls{6xhis}-\gls{utp_E} expressed by EC174, which was grown in \gls{2yt} at 28°C and harvested 4 hours after induction as its expression condition.}
       \label{fig:TestPuri_WB}
\end{figure}

The test purification was done with the supernatant sample from the condition containing the highest amount of soluble protein. The Western blot of both \gls{6xhis}-tagged \gls{utp} showed that they required similar concentration of imidazole, in the range of 250~--~500~mM for elution, with \gls{utp_E} eluting earlier than \gls{utp_C} (see Figures~\ref{fig:TestPuri_WB}). Though when these bands were compared to the \glsxtrlong{ft} and wash fraction, it became apparent that the \gls{utp}s rather remained unbound to the \gls{Ni-NTA} resin. Moreover, a smaller protein (around 30~kDa) eluted at around 50~--~100 mM imidazole in both uTP samples, correlating with its previous occurrences in the Western blots of \gls{utp_C} (see Supplementary Figures~\ref{sup:Test_WB_C}~and~\ref{sup:SolTest_Western}B) 


%---------------------------------------------------------------------------------------------------
%%%%%%%%%%%%%%%%%%%%%%%%%%%%%%%%%%%%%%%%%%%%%%%%%%%%%%%%%%%%%%%%%%%%%%%%%%%%%%%%%%%
%---------------------------------------------------------------------------------------------------
\section{Bioinformatic analysis}
The \gls{utp}s were also further characterized through bioinformatic analysis, such as structure prediction, alignment, homology searches, and other analyses using the tools listed in Table~\ref{tab:software_list}. Within the scope of this thesis, the analyses were limited to the three \gls{utp}s (\gls{utp_B}, \gls{utp_E}, \gls{utp_C}). 
%---------------------------------------------------------------------------------------------------
\subsection{Hypothetical structures of the uTPs}
Every  \gls{utp}'s structure was predicted by Alphafold 3. The predicted structures had two parts, an N- and C-terminal region (see Figure~\ref{fig:Alphafold_structure}). The C-terminal region consisted of disordered stretches as well as α-helices, though there were slight differences between the three \gls{utp}s in the number of α-helices present and in their localization. However, the prediction quality for this region was low. In contrast, the model was more confident in the structure of the N-terminal region. An alignment of the first 90 amino acids was done using Pymol (see Figure~\ref{fig:Aln_structure}), which showed that the N-terminal region was predicted to contain multiple α-helices that were tightly packed and arranged in a U-shape. Furthermore, the \gls{utp}s were also predicted with their respective cargo protein and \gls{egfp} and behaved the same as predicted individually. 

\begin{figure}[H]
        \centering
        \includegraphics[width=16cm]{"4.3.1_Alphafold_structure"}
\caption{
\textbf{Structure Prediction of three \gls{utp}s by Alphafold3}. 
    Structures were predicted using the Alphafold 3 server  (\cite{abramson_accurate_2024}; \url{https://alphafoldserver.com}). The predicted structures were colored by the confidence level (pLDDT score) of the model; marked termini (N, C).}
       \label{fig:Alphafold_structure}
\end{figure}

\begin{figure}[H]
        \centering
        \includegraphics[width=7.5cm]{"4.3.1_Structure_Aln"}
\caption{
\textbf{Structural alignment of the first 90 amino acids of the predicted uTP structures}. 
    The predicted uTP structures (see Figure~\ref{fig:Alphafold_structure}) were aligned using Pymol (Schrödinger). Alignment was done for the first 90 amino acids (N-terminal region of the uTPs). The colors represented: blue for \gls{utp_B}, red for \gls{utp_E}, and grey for \gls{utp_C}.}
       \label{fig:Aln_structure}
\end{figure}

%---------------------------------------------------------------------------------------------------
%---------------------------------------------------------------------------------------------------
\subsection{Alignment and physicochemical properties}
An alignment of the three \gls{utp}s was done with Clustal Ω (see Figure~\ref{fig:aln_omega}). The uTPs only shared around 30 to 40~\% identity but still contained conserved regions. These regions were predominantly hydrophobic and correlated with the prediction of α-helices by Alphafold. Furthermore, Polar residues were interspersed within and between the conserved regions.

\begin{figure}[H]
        \centering
        \includegraphics[width=13cm]{"4.3.2_omega_aln"}
\caption{
\textbf{Clustal Ω alignment of the uTP sequences studied in this work}. 
    Alignment of the three \gls{utp}s (\gls{utp_B}, \gls{utp_E}, \gls{utp_C}) was done with Clustal Ω (\cite{sievers_fast_2011}; \url{https://www.ebi.ac.uk/jdispatcher/msa/clustalo}); colored by conserved regions and visualized with JalView.}
       \label{fig:aln_omega}
\end{figure} 

Additionally, these sequences were analyzed with HeliQuest, which returned values for standard hydrophobicity and the hydrophobic moment, a score for how polar and non-polar amino acids were distributed in an α-helix \parencite{gautier_heliquest_2008}. These values were then plotted against the starting residue of the window that it calculated the helix properties for (see Supplementary Figure\ref{sup:hydro_plots}). The mean hydrophobic moment was around 0.2, indicating amphipathic character for the putative helices. Moreover, the first half of the protein contained more predicted α-helices that displayed an amphipathic nature, whereas α-helices that form in the latter half of the peptide sequence were distinctively more hydrophobic than amphipathic. Furthermore, the amphipathic α-helices tended to have a polar face without a specific charge accompanying them (see exemplary plots in Supplementary Figure~\ref{sup:helices}).


%---------------------------------------------------------------------------------------------------
%---------------------------------------------------------------------------------------------------
\subsection{Structural homology searches and signal prediction}
A structural homology search was done with PHYRE \parencite{powell_phyre22_2025}. For that purpose, the uTPs were segmented into the N-terminal (first 90 amino acids) and C-terminal (residues 91~--~180) region. Hits obtained by PHYRE had overall low to medium confidence and low coverage. Both regions were associated with α-helical structures that were also transmembrane domains. The highest confidence a model obtained was 82~\% for the C-terminal region, specifically residues 140~--~169, of \gls{utp_E} with the model being the A-chain of ERBB4 intracellular domain (PF21314), a transmembrane domain. 

\par The sequences were further analyzed with DeepLoc \parencite{odum_deeploc_2024}. DeepLoc predicted the \gls{utp}s to have membrane associations with transmembrane domains, though it did not predict the localization. Furthermore,  this tool postulated that the residues 140~--~170 were most important for signal sorting, while the last few residues contributed slightly. Only for the \gls{utp_C} there was an additional region, amino acids 124~--~134, that contributed to the importance of the signal sorting.
%---------------------------------------------------------------------------------------------------
%%%%%%%%%%%%%%%%%%%%%%%%%%%%%%%%%%%%%%%%%%%%%%%%%%%%%%%%%%%%%%%%%%%%%%%%%%%%%%%%%%%
